\documentclass[14pt,dvipdfmx]{jsarticle}
\input{lecture-preamble.tex}

\begin{document}

\lectureheader[4]{数学Ⅱ}{6}{微分法と積分法}{1}{微分係数と導関数}{2-B}{関数の微分}
  {184,185}{定数倍と和・差の性質を使って, 整式を微分しよう}

% 1-2-A の定理 6.1.1 の続き（定義は 6.1.5 まで）
\setcounter{theorem}{1}

% ============ 1ページ目（表・左） ============
\vspace{-2.0zw}%
\heading{〜前時の復習〜}
\vspace{-0.5em}
\begin{reviewbox}{導関数}
  \begin{enumerate}[label=\textbf{\arabic*}, leftmargin=1.5zw, itemsep=0.08em]
    \item 導関数は
          \par\vspace{0.15em}
          {\renewcommand{\fitblankpad}{1.2zw}%
          \fitblankbf{f'(x)=\lim_{h \to 0}\frac{f(x+h)-f(x)}{h}}}
          \mbox{である.}
          \par\vspace{0.25em}
    \item $(x^{n})'=\fitblankbf{nx^{n-1}}$,\quad
          $(c)'=\fitblankbf{0}$
          \quad（$n$ は正の整数, $c$ は定数）
  \end{enumerate}
\end{reviewbox}

\vspace{0.15em}
\heading{関数の微分}
\vspace{-0.45em}
関数 $f(x)$ から導関数 $f'(x)$ を求めることを,
$f(x)$ を $x$ で\fitblankbf{\text{微分する}}という.
単に微分するともいう.
\par\vspace{0.2em}
\textbf{注意}\quad
今後, 関数 $f(x)$, $g(x)$ は $x$ の多項式で表される関数とする.
\par\vspace{0.3em}
2つの関数 $f(x)=x^{3}$, $g(x)=x^{2}$ の導関数
{\renewcommand{\fitblankpad}{0.7zw}%
\mbox{$f'(x)=\fitblankbf{3x^{2}}$,}\ \mbox{$g'(x)=\fitblankbf{2x}$}}
をもとに, 関数 $y=2f(x)$ や $y=f(x)+g(x)$ などを微分してみよう.
\par\vspace{0.25em}
まず, 関数 $y=2f(x)$ すなわち $y=2x^{3}$ の導関数を求めよう.
\vspace{1em}
\begin{answerbox}{}[15]
  $y'=\lim_{h \to 0}\dfrac{2(x+h)^{3}-2x^{3}}{h}
    =\lim_{h \to 0}2\left\{\dfrac{(x+h)^{3}-x^{3}}{h}\right\}$\\[0.4em]
  $=\lim_{h \to 0}2(3x^{2}+3xh+h^{2})
    =\lim_{h \to 0}(2\cdot 3x^{2}+2\cdot 3xh+2h^{2})$\\[0.4em]
  $=2\cdot 3x^{2}$
  \boxanswer{$2\cdot 3x^{2}$}
\end{answerbox}

\vfill
\nextpage

% ============ 2ページ目（表・右） ============
\vspace{-1.4zw}%
すなわち $y'=\fitblankbf{2\cdot 3x^{2}}$
\par\vspace{0.35em}
ここで, $f'(x)=\fitblankbf{3x^{2}}$ であるから, 次のことがいえる.
\par\vspace{0.3em}
$y=2f(x)$ を微分すると \quad $y'=\fitblankbf{2f'(x)}$

\vspace{1.5em}
次に, 関数 $y=f(x)+g(x)$ すなわち $y=x^{3}+x^{2}$ の導関数を求めよう.

\begin{answerbox}{}[15]
  $y'=\lim_{h \to 0}\dfrac{\{(x+h)^{3}+(x+h)^{2}\}-(x^{3}+x^{2})}{h}$\\[0.45em]
  $=\lim_{h \to 0}\left\{\dfrac{(x+h)^{3}-x^{3}}{h}+\dfrac{(x+h)^{2}-x^{2}}{h}\right\}$\\[0.45em]
  $=\lim_{h \to 0}\{(3x^{2}+3xh+h^{2})+(2x+h)\}$\\[0.4em]
  $=3x^{2}+2x$
  \boxanswer{$3x^{2}+2x$}
\end{answerbox}
\vspace{0.5em}
すなわち $y'=\fitblankbf{3x^{2}+2x}$
\par\vspace{0.35em}
ここで, $f'(x)=\fitblankbf{3x^{2}}$, $g'(x)=\fitblankbf{2x}$ であるから,
次のことがいえる.
\par\vspace{0.3em}
$y=f(x)+g(x)$ を微分すると \quad $y'=\fitblankbf{f'(x)+g'(x)}$
\par\vspace{0.55em}
同様にして, 次のこともいえる.
\par\vspace{0.3em}
$y=f(x)-g(x)$ を微分すると \quad $y'=\fitblankbf{f'(x)-g'(x)}$
\vspace{1.5em}

一般に, 関数 $f(x)$, $g(x)$ について, 次の性質が成り立つ.

\vspace{-0.15em}
\begin{theorembox}{関数の定数倍および和, 差の導関数}
  $k$ を定数とする.
  \par\vspace{0.2em}
  \begin{enumerate}[label=\textbf{\arabic*}, leftmargin=1.6zw, itemsep=0.45em]
    \item $y=kf(x)$ を微分すると \quad
          $y'=\fitblankbf{kf'(x)}$
    \item $y=f(x)+g(x)$ を微分すると \quad
          $y'=\fitblankbf{f'(x)+g'(x)}$
    \item $y=f(x)-g(x)$ を微分すると \quad
          $y'=\fitblankbf{f'(x)-g'(x)}$
  \end{enumerate}
\end{theorembox}

\vfill
\nextpage

% ============ 3ページ目（裏・左） ============
\vspace{-1.4zw}%
\noindent{\bfseries 例\,6}\quad
定理 6.1.2 を用いて, 次の関数を微分する.
\begin{tasks}[label=(\arabic*), label-width=1.6zw, item-indent=2zw, after-item-skip=0.08em](2)
  \task 関数 $y=3x^{2}-4x+2$
  \task 関数 $y=\dfrac{2}{3}x^{3}-\dfrac{5}{2}x^{2}+5x$
\end{tasks}
\vspace{-0.55em}
\noindent
\begin{minipage}[t]{0.49\linewidth}
\begin{answerbox}{(1)}[16]
  $y'=(3x^{2})'-(4x)'+(2)'$\\[0.2em]
  $=3(x^{2})'-4(x)'+(2)'$\\[0.2em]
  $=3\cdot 2x-4\cdot 1+0$\\[0.2em]
  $=6x-4$
  \boxanswer{$6x-4$}
\end{answerbox}
\end{minipage}\hfill
\begin{minipage}[t]{0.49\linewidth}
\begin{answerbox}{(2)}[16]
  $y'=\dfrac{2}{3}(x^{3})'-\dfrac{5}{2}(x^{2})'+5(x)'$\\[0.25em]
  $=\dfrac{2}{3}\cdot 3x^{2}-\dfrac{5}{2}\cdot 2x+5\cdot 1$\\[0.25em]
  $=2x^{2}-5x+5$
  \boxanswer{$2x^{2}-5x+5$}
\end{answerbox}
\end{minipage}

\vspace{0.15em}
\begin{practiceprob}{8}
  定理 6.1.2 を用いて, 次の関数を微分せよ.
  \begin{tasks}[label=(\arabic*), label-width=1.6zw, item-indent=2zw, after-item-skip=0.06em](2)
    \task $y=4x^{2}+3x-4$
    \task $y=2x^{2}-5x+1$
    \task $y=-3x^{2}+x-2$
    \task $y=-x^{2}-x+3$
  \end{tasks}
\end{practiceprob}
\vspace{-0.7em}
\noindent
\begin{minipage}[t]{0.49\linewidth}
\begin{answerbox}{(1)}[12]
  $y'=4(x^{2})'+3(x)'+(-4)'$\\[0.25em]
  $=4\cdot 2x+3\cdot 1+0=8x+3$
\end{answerbox}
\end{minipage}\hfill
\begin{minipage}[t]{0.49\linewidth}
\begin{answerbox}{(2)}[12]
  $y'=2(x^{2})'-5(x)'+(1)'$\\[0.25em]
  $=2\cdot 2x-5\cdot 1+0=4x-5$
\end{answerbox}
\end{minipage}

\noindent
\begin{minipage}[t]{0.49\linewidth}
\begin{answerbox}{(3)}[12]
  $y'=-3(x^{2})'+(x)'+(-2)'$\\[0.25em]
  $=-3\cdot 2x+1+0=-6x+1$
\end{answerbox}
\end{minipage}\hfill
\begin{minipage}[t]{0.49\linewidth}
\begin{answerbox}{(4)}[12]
  $y'=-(x^{2})'-(x)'+(3)'$\\[0.25em]
  $=-2x-1+0=-2x-1$
\end{answerbox}
\end{minipage}

\vfill
\nextpage

% ============ 4ページ目（裏・右） ============
\vspace{-1.4zw}%
\noindent{\bfseries 練習\,8}\quad（続き）
\begin{tasks}[label=(\arabic*), label-width=1.6zw, item-indent=2zw, after-item-skip=0.08em, start=5](2)
  \task $y=x^{3}+2x^{2}-3x$
  \task $y=-2x^{3}-x^{2}+6x-2$
  \task $y=\dfrac{4}{3}x^{3}+\dfrac{3}{4}x^{2}-\dfrac{1}{2}x$
  \task $y=-\dfrac{1}{3}x^{3}-\dfrac{3}{2}x-\dfrac{1}{2}$
\end{tasks}
\vspace{0.2em}
\noindent
\begin{minipage}[t]{0.49\linewidth}
\begin{answerbox}{(5)}[10]
  $y'=(x^{3})'+2(x^{2})'-3(x)'$\\[0.3em]
  $=3x^{2}+4x-3$
\end{answerbox}
\end{minipage}\hfill
\begin{minipage}[t]{0.49\linewidth}
\begin{answerbox}{(6)}[10]
  $y'=-2(x^{3})'-(x^{2})'+6(x)'$\\[0.3em]
  $=-6x^{2}-2x+6$
\end{answerbox}
\end{minipage}

\vspace{0.35em}
\begin{answerbox}{(7)}[11]
  $y'=\dfrac{4}{3}(x^{3})'+\dfrac{3}{4}(x^{2})'-\dfrac{1}{2}(x)'$\\[0.35em]
  $=\dfrac{4}{3}\cdot 3x^{2}+\dfrac{3}{4}\cdot 2x-\dfrac{1}{2}$\\[0.3em]
  $=4x^{2}+\dfrac{3}{2}x-\dfrac{1}{2}$
\end{answerbox}
\vspace{0.15em}
\begin{answerbox}{(8)}[11]
  $y'=-\dfrac{1}{3}(x^{3})'-\dfrac{3}{2}(x)'+\left(-\dfrac{1}{2}\right)'$\\[0.35em]
  $=-\dfrac{1}{3}\cdot 3x^{2}-\dfrac{3}{2}$\\[0.3em]
  $=-x^{2}-\dfrac{3}{2}$
\end{answerbox}

\vfill
\begin{summarybox}
  \begin{enumerate}[label=\textbf{\arabic*}, leftmargin=1.5zw, itemsep=0.2em, topsep=0.15em]
    \item $f(x)$ を微分するとは, 導関数 $f'(x)$ を求めることである.
    \item $k$ を定数とする.\enspace
          {\renewcommand{\fitblankpad}{0.6zw}%
          $(kf(x))'=\fitblankbf{kf'(x)}$,\
          $(f(x)+g(x))'=\fitblankbf{f'(x)+g'(x)}$,\
          $(f(x)-g(x))'=\fitblankbf{f'(x)-g'(x)}$}
  \end{enumerate}
\end{summarybox}

\vspace{0.05em}
\nexttime{積は展開してから微分し, 微分係数を導関数で求めてみよう!}

\end{document}
